\documentclass[11pt,letterpaper]{article}
\usepackage[T1]{fontenc}
\usepackage[utf8]{inputenc}
\usepackage{lmodern,microtype}
\usepackage[margin=1in,headheight=15pt]{geometry}
\usepackage{amsmath,amssymb,amsthm,mathtools}
\usepackage{booktabs,longtable,array,enumitem,xcolor,fancyhdr,listings}
\usepackage[hidelinks]{hyperref}
\usepackage{bookmark}
\definecolor{ink}{HTML}{183246}
\definecolor{accent}{HTML}{216778}
\definecolor{pale}{HTML}{F1F5F7}
\pagestyle{fancy}
\fancyhf{}
\fancyhead[L]{\small Template-free equality in iPRA}
\fancyhead[R]{\small Research note, November 2026}
\fancyfoot[C]{\thepage}
\setlength{\parskip}{4pt}
\setlength{\parindent}{0pt}
\setlist{itemsep=3pt,topsep=5pt}
\emergencystretch=2em
\newtheorem{theorem}{Theorem}[section]
\newtheorem{lemma}[theorem]{Lemma}
\newtheorem{proposition}[theorem]{Proposition}
\newtheorem{corollary}[theorem]{Corollary}
\theoremstyle{definition}
\newtheorem{definition}[theorem]{Definition}
\theoremstyle{remark}
\newtheorem{remark}[theorem]{Remark}
\newcommand{\D}{\mathsf D}
\newcommand{\Q}{\mathsf Q}
\newcommand{\Ncal}{\mathcal N}
\newcommand{\Gcal}{\mathcal G}
\newcommand{\ipra}{\mathsf{iPRA}}
\newcommand{\fv}{\operatorname{FV}}
\newcommand{\Rep}[3]{R_{#1}^{#2}(#3)}
\newcommand{\code}[1]{\texttt{\detokenize{#1}}}
\lstset{basicstyle=\ttfamily\footnotesize,breaklines=true,columns=fullflexible,
 keepspaces=true,showstringspaces=false,frame=single,rulecolor=\color{accent},
 backgroundcolor=\color{pale},aboveskip=10pt,belowskip=10pt}
\hypersetup{pdftitle={Template-free equality in a seven-rule presentation of iPRA},
 pdfauthor={Research prepared for Vladimir Reshetnikov},
 pdfsubject={Exact sequent equivalence, bi-interpretability and size audit}}
\title{\color{ink}\bfseries Template-free equality in a\newline seven-rule presentation of iPRA}
\author{Research prepared for Vladimir Reshetnikov}
\date{November 2026}
\begin{document}
\maketitle
\begin{abstract}
The supplied equality template can be eliminated from the native
ConciseIPRA presentation. Replace equality transport by the deterministic
rule
\[
 \frac{s=t\qquad u=v}{v=R_s^t(u)},
\]
where $R_s^t$ replaces every exact free occurrence of $s$ by $t$ in a
single root-first pass. The rule takes only two premise proofs: no
template, path, mask, or selected occurrence is supplied.

The main lemma recovers selective replacement from this global operation.
A computation term chosen larger than the entire desired context marks
exactly the occurrences to be changed. Computation at zero and the new
rule alone then recover all ordinary equational transport. Combined with
the repository's existing witness-sensitive introduction, the resulting
seven-rule calculus proves exactly the native sequents, in arbitrary
contexts. Its full first-order theory is consequently bi-interpretable
and equiconsistent with iPRA, with the same numerical domain and actual
equality.

This is a precise reduction of the primitive equality interface, not a
proof of absolute minimality. Complete checkers, matched comparisons,
proof-producing examples, rejection controls and reproducible measurements
are included. The tagged checker improves on the comparable one-template
implementation. It does not beat the smallest compact programs, whose
different input contracts are explicitly distinguished.
\end{abstract}

\clearpage
\tableofcontents
\clearpage

\section{The result and its comparison boundary}
\label{new:status}

The proposal is a calculus $\Gcal$ with the same numerical terms and
native formulas as the repository's $\Ncal$ (\code{NPrf}). Its central
property is
\[
             \Gamma\vdash_{\Gcal}A
             \quad\Longleftrightarrow\quad
             \Gamma\vdash_{\Ncal}A
             \tag{1}\label{new:exact}
\]
for every finite context and every native formula, including open
sequents. Thus the change does not merely preserve closed equations or
consistency strength. It preserves the formula language and all
derivability needed for both interpretation composites.

The witness-sensitive introduction is inherited from
\code{NativeOneIntro.lean}~\cite{native}. The contribution here is to replace a
user-supplied equality context by a fixed global syntax operation, and to
prove that no equality strength is lost. The completeness argument uses
an oversized computation marker; it does not invoke the logical rules,
induction, a numerical case distinction, or substitution into a derivation.

\subsection{The precise simplification}

Native equality transport supplies an equation template, or equivalently
two numerical templates. The newer \code{one_context} proposal supplies
one numerical template. The present rule supplies none:
\[
\begin{array}{c@{\qquad}c}
\text{one-context predecessor}&\text{this proposal}\\[2pt]
\displaystyle\frac{s=t\quad H(s)=v}{v=H(t)}
&
\displaystyle\frac{s=t\quad u=v}{v=R_s^t(u)}.
\end{array}
\]
The selected occurrences become a derived proof construction, rather than
a parameter of an individual equality inference. This is the exact
mathematical simplification claimed.

It is not enough to count rule headings: arbitrary machinery could be
hidden in a rule. We therefore define $R$ completely, provide an executable
checker, and measure whole programs. A smaller primitive certificate
interface need not give the smallest implementation or the shortest
proofs. In particular, the current compact byte record is not improved
by this implementation. No language-independent simplest possible
axiomatization is established.

\subsection{Pinned sources and the moving proposal}

The inherited mathematical sources were copied from ProveIt commit
\begin{center}\small\ttfamily
a4e5e96c456303cf231eab2aeeb3291d8889833f.
\end{center}
The source directory is
\code{Logic/PrimitiveRecursiveArithmetic/ConciseIPRA}.
Paths, immutable URLs, Git blob identities and SHA-256 hashes are in
\code{evidence/sources.json}. The copied files and original MIT No
Attribution license are included.

The repository advanced during the investigation. The filed comparison
was updated through \code{521884df22e8f2dac0b7066095a42ae975ea4166},
then checked once more at the frozen final boundary
\begin{center}\small\ttfamily
2a724ea8ec6665f91aa74c892414584b36ca698d.
\end{center}
The latter check includes the additional \code{checked_refinement},
\code{one_sided_transport} and \code{validated_core} archives.

In particular, \code{one_context} already contains the one-template
crossed rule above. That rule is acknowledged as a predecessor, not
claimed as a new result here. The \code{kernel_refinement} package also
contains a supplied single-occurrence context rule~\cite{comparison}. Comparison sources,
hashes and the bounded novelty audit are recorded in
\code{evidence/comparison_sources.json} and
\code{evidence/comparison_boundary.json}.

The novelty claim is limited to the inspected proposal and packages.
Equality by rewriting is familiar mathematics; no literature-wide
priority claim is made.

\section{The complete primitive presentation}
\label{new:rules}

\subsection{Numerical terms, formulas and scopes}

There is one numerical sort. Terms and formulas are finite expressions
\[
 t ::= x\mid 0\mid S(t)\mid I(n,a;u.h),
 \qquad
 A ::= (s=t)\mid \Q_y A\mid \D_{x,y}A .
\]
In $I(n,a;u.h)$, only $u$ is bound, and only in $h$; the counter and seed
are outside its scope. The update may contain fixed numerical parameters.
The two variables bound by $\D$ are distinct.

The intended ordinary readings are
\begin{align*}
 \Q_y B(y)&:=B(0)\to\forall y\,B(y),\\
 \D_{x,y}B(x,y)&:=\exists x\bigl(B(x,0)\to\forall y\,B(x,y)\bigr).
\end{align*}
These are interpretations into ordinary intuitionistic logic, not extra
primitive axioms. The designated contradiction is $0=S0$. All substitutions
are capture avoiding; bound variables are identified up to alpha
conversion, or represented canonically by de Bruijn indices.

Contexts $\Gamma$ are finite reusable collections of formulas. Proofs are
finite trees. Computation axioms are available in every context.

\subsection{The fixed replacement operation}
\label{new:replacement}

For terms $s,t,u$ in the same numerical context, define $\Rep{s}{t}{u}$
by a single structural traversal:
\begin{enumerate}
\item Test the current subtree for exact equality with the appropriately
lifted copy of $s$. If it matches, return the corresponding lifted copy
of $t$, and stop traversing this subtree.
\item Otherwise preserve the constructor and recurse on its term
children. Under an iterator's update binder, lift both $s$ and $t$ by
one additional binding level.
\item An unmatched variable or zero is unchanged.
\end{enumerate}
The root test precedes recursion. Inserted targets are never traversed,
and newly created matches are never revisited. Thus replacement is
total on finite syntax even when $t$ contains $s$. In particular,
\[
                    \Rep{s}{t}{s}=t.
\]
A bound state variable must not be confused with a free source variable.
For example, replacing a free $x$ by $0$ under $u.h$ changes lifted
occurrences of $x$ and leaves the bound $u$ untouched.

This is deterministic replacement of all exact free occurrences in one
pass. It is not repeated rewriting to a normal form. No termination
ordering, normalization theorem or unbounded search is part of the rule.

\subsection{Exactly seven rules}

\paragraph{1. Hypothesis.}
\[
 \frac{A\in\Gamma}{\Gamma\vdash A}. \tag{Hyp}
\]

\paragraph{2--3. Computation.}
For every well-scoped instance,
\begin{align*}
 \Gamma&\vdash I(0,a;u.h)=a, &&\tag{$I_0$}\\
 \Gamma&\vdash I(Sn,a;u.h)=h[I(n,a;u.h)/u]. &&\tag{$I_S$}
\end{align*}

\paragraph{4. Global crossed replacement.}
\[
 \frac{\Gamma\vdash s=t\qquad\Gamma\vdash u=v}
      {\Gamma\vdash v=\Rep{s}{t}{u}}. \tag{G}\label{new:grule}
\]
All four endpoints are numerical terms. Sources are unrestricted,
including literal zero. There is no supplied term or formula template.
Reflexivity, symmetry, transitivity and general congruence are not
additional primitive rules.

\paragraph{5. Closure elimination.}
\[
 \frac{\Gamma\vdash\Q_yB(y)\qquad\Gamma\vdash B(0)}
      {\Gamma\vdash B(t)}. \tag{QE}
\]
The term $t$ is arbitrary, subject to scope and capture avoidance.

\paragraph{6. Paired-binder elimination.}
\[
 \frac{\Gamma\vdash\D_{x,y}B(x,y)\qquad
       \Gamma,\Q_yB(b,y)\vdash A}
      {\Gamma\vdash A}. \tag{DE}
\]
The eigenparameter $b$ is fresh for $\Gamma,A$ and the free parameters of
$B$. It is local to the second premise.

\paragraph{7. Witness-sensitive introduction.}
Define a decidable test of finite syntax by
\[
 \eta(B,t,z)=
 \begin{cases}
 z,&\text{$B$ is a literal equation and $t$ is literally $0$},\\
 0,&\text{otherwise}.
 \end{cases}
\]
Then the single introduction rule is
\[
 \frac{\Gamma,B(t,\eta(B,t,z))\vdash B(t,Sz)}
      {\Gamma\vdash\D_{x,y}B(x,y)}. \tag{DI}
\]
Here $z$ is fresh for $\Gamma,t$ and the free parameters of $B$.
Literal zero means the zero constructor, not a term provably equal to
zero. The witness-sensitive test is inherited from the existing proposal.

There are no further logical or arithmetic axioms or inference rules.
Successor injectivity, separation, numerical coverage and induction are
derived through the equivalence below. Structural rules are metatheorems,
not extra proof packets.

\section{Recovering all equality reasoning}
\label{new:equality}

The following argument uses only $I_0$ and G. It is valid in every
retained context, with arbitrary numerical parameters.

\subsection{The elementary equality operations}

\begin{lemma}[Common-left joining]\label{new:join}
From $s=t$ and $s=u$, derive $u=t$.
\end{lemma}
\begin{proof}
Use the two equations as the major and minor premises of G. Its
replacement starts with a root match, so $\Rep{s}{t}{s}=t$.
\end{proof}

\begin{lemma}[Reflexivity, symmetry and transitivity]\label{new:algebra}
The usual elementary equality operations are derivable.
\end{lemma}
\begin{proof}
Put $J(a)=I(0,a;w.0)$. The axiom $J(a)=a$, used twice in
Lemma~\ref{new:join}, yields $a=a$.

Given $a=b$, join the major $a=a$ with the minor $a=b$ to obtain $b=a$.

Given $a=b$ and $b=c$, first obtain $b=a$ by symmetry; join the major
$b=c$ with the minor $b=a$ to obtain $a=c$.
\end{proof}

No general congruence is assumed in these derivations. Only a root
match is used.

\subsection{A marker that cannot occur accidentally}

Let $H(x)$ be a numerical term template with a distinguished free variable
$x$, distinct from the passive parameters in $s,t$. There may be any
finite number of distinguished occurrences, including zero.
Write $|T|$ for the number of constructor and variable nodes of a term,
including the update body of every iterator. Renaming and de Bruijn
lifting preserve this number.

Choose
\[
 N=1+|H|+|s|+|t|,\qquad
 K=I(0,s;w.S^N0).
 \tag{2}\label{new:marker}
\]
The closed update ignores $w$. The whole term $K$ need not be closed:
it retains the free parameters of $s$. Its syntax is larger than $H,s,t$,
and $I_0$ immediately proves $K=s$.

\begin{lemma}[Exact marking]\label{new:markerlemma}
The only capture-safe occurrences of $K$ in $H(K)$ are the copies
inserted at the distinguished occurrences of $x$. Consequently,
\[
 \Rep{K}{s}{H(K)}=H(s),\qquad
 \Rep{K}{t}{H(K)}=H(t)
 \tag{3}\label{new:markeridentities}
\]
as finite syntax, up to the agreed binder convention.
\end{lemma}
\begin{proof}
Consider a possible occurrence of a subtree equal to the appropriately
lifted $K$. Tree occurrences are disjoint or nested.

If it is disjoint from every inserted copy of $K$, it is entirely inside
a passive part of $H$, which has fewer nodes than $K$.
If it is properly inside an inserted copy, it is a proper subtree of
that copy and again has fewer nodes.
If it properly contains an inserted copy, it has more nodes than $K$.
None of these can be equal to $K$.

The remaining possibility is an entire inserted copy of $K$. Under a
binder the planted copy is lifted, and the comparison source is lifted
by the same amount. Lifting preserves node count, so the argument is
unchanged. G stops at each matching root and replaces exactly those
copies, which gives both displayed syntactic identities.
\end{proof}

The identities in (3) describe a finite syntax algorithm. They are not
object-theoretic uses of the equality rule we are trying to derive.

\subsection{Selective congruence from global replacement}

\begin{theorem}[Arbitrary numerical congruence]\label{new:congruence}
For every numerical template $H$,
\[
           \Gamma\vdash s=t
           \quad\Longrightarrow\quad
           \Gamma\vdash H(s)=H(t).
\]
\end{theorem}
\begin{proof}
Choose $K$ as in (2). The axiom $K=s$ and the premise $s=t$ give $K=t$
by derived transitivity. Derived reflexivity gives $H(K)=H(K)$.

Apply G with major $K=s$ and this reflexive minor. By (3), the result is
$H(K)=H(s)$. Apply G with major $K=t$ and the same reflexive minor.
The result is $H(K)=H(t)$. Common-left joining, with the latter equation
as major, gives $H(s)=H(t)$.
\end{proof}

The compiler uses $H$ to construct ordinary terms and computation
certificates. The primitive G packet still contains only two proofs.
The checker receives no hidden occurrence set, template or proof of a
syntactic side condition. Certificate size may increase.

\begin{theorem}[Full equational transport]\label{new:transport}
The following rule is derivable:
\[
 \frac{\Gamma\vdash s=t\qquad\Gamma\vdash a(s)=b(s)}
      {\Gamma\vdash a(t)=b(t)}.
 \tag{Tr}
\]
\end{theorem}
\begin{proof}
Theorem~\ref{new:congruence} gives $a(s)=a(t)$ and $b(s)=b(t)$.
Derived symmetry and transitivity compose the chain
\[
                  a(t)=a(s)=b(s)=b(t).
\]
\end{proof}

\paragraph{Why passive occurrences matter.}
Suppose $s=t$ and $H(x)=I(x,s;w.s)$, where the two displayed passive
copies of $s$ are to stay fixed. Directly applying $R_s^t$ to $H(s)$
would also replace them. Instead form $H(K)=I(K,s;w.s)$.
Since $K$ is too large to occur in the passive parts, the two global
replacements produce exactly
\[
                  I(s,s;w.s),\qquad I(t,s;w.s).
\]
The marker proof therefore establishes selective congruence even in
the case where naive all-occurrence replacement is insufficient.

\subsection{The converse: G is native equality reasoning}

\begin{lemma}\label{new:gsound}
Every G inference is derivable in the native calculus.
\end{lemma}
\begin{proof}
Mark the occurrences that the fixed root-first algorithm replaces in
$u$. They are disjoint because traversal stops at a match. Abstract just
those occurrences by a fresh hole to obtain a numerical template $H$.
Capture-safe matching ensures, as syntax,
\[
                       u=H(s),\qquad R_s^t(u)=H(t).
\]
Native equational transport changes $H(s)=v$ to $H(t)=v$.
Native symmetry gives $v=H(t)$, exactly G's conclusion. If no occurrence
is replaced, $H$ is constant, which is permitted.
\end{proof}

Thus $I_0$ together with G recovers all the equality machinery that G
itself needs for soundness in the old calculus. This equality bootstrap
has no dependence on induction or the introduction rules.

\section{Recovering the native logical rules}
\label{new:intro}

\subsection{Structural preliminaries}

\begin{lemma}[Renaming, weakening and cut]\label{new:struct}
Injective renaming of free parameters, alpha renaming of local parameters,
weakening of retained contexts and cut are admissible in $\Gcal$.
\end{lemma}
\begin{proof}
Injective renaming preserves exact syntax matches and commutes with the
capture-safe replacement operation. It also preserves the equation-root
and literal-zero tests in DI.

For weakening, induct on a derivation and rename a local eigenparameter
away from added assumptions. For cut, rename local eigenparameters away
from the finite data of the inserted derivation, then replace uses of
the cut assumption by that derivation. The original numerical terms,
matrix roots and witness syntax are unchanged.
\end{proof}

Arbitrary numerical substitution is deliberately not asserted by a direct
induction here. It can create new exact matches for G, or turn a variable
witness into literal zero and change the DI branch. Its admissibility
will follow from exact native equivalence.

\subsection{Recovering unrestricted case introduction}

The old native case rule is
\[
 \frac{\Gamma,B(t,0)\vdash B(t,Sz)}
      {\Gamma\vdash\D_{x,y}B(x,y)}. \tag{DI-case}
\]

\begin{lemma}\label{new:dcase}
Every DI-case instance is derivable in $\Gcal$.
\end{lemma}
\begin{proof}
Primitive DI already gives the result unless $B$ is an equation and $t$
is literally zero. In that exception put $w=I(0,0;u.0)$.
This is syntactically an iterator, and $I_0$ proves $w=0$.

Under $B(w,0)$, derived equational transport gives $B(0,0)$.
Cut in the supplied case branch to obtain $B(0,Sz)$, and transport along
$0=w$ to obtain $B(w,Sz)$. These transports are equational because the
exceptional matrix is an equation.

Apply primitive DI at witness $w$. Its syntax selects the case branch
and concludes the original $\D_{x,y}B(x,y)$. No witness term occurs
freely in that conclusion. Choose the local variables fresh.
\end{proof}

This step is not circular: full equational transport was established
from $I_0$ and G alone.

\subsection{Recovering direct equation progress}

The native progress rule is
\[
 \frac{\Gamma,E(z)\vdash E(Sz)}
      {\Gamma\vdash\Q_yE(y)},\qquad E\text{ a literal equation}.
 \tag{QI-progress}
\]

\begin{lemma}\label{new:qprog}
Every QI-progress instance is derivable in $\Gcal$.
\end{lemma}
\begin{proof}
Introduce a dummy witness coordinate $x$ absent from $E$.
Primitive DI with witness zero, matrix $B(x,y)=E(y)$, and the supplied
progressive premise gives $\D_{x,y}E(y)$. Apply DE.
Its local assumption is $\Q_yE(y)$ and is independent of the fresh
witness. Hyp therefore closes the branch with that same formula as
conclusion, and the eigenvariable condition holds.
\end{proof}

Dummy-coordinate DI-case similarly derives case introduction of $\Q$
for arbitrary native matrices. Equation induction is obtained from
QI-progress and QE after supplying its base case.

\section{Exact equivalence and both interpretations}
\label{new:equivalence}

\begin{theorem}[Exact sequent equivalence]\label{new:main}
For every finite native context $\Gamma$ and formula $A$,
\[
       \Gamma\vdash_{\Gcal}A
       \quad\Longleftrightarrow\quad
       \Gamma\vdash_{\Ncal}A.
\]
Both transformations of finite proofs are effective. Numerical terms,
formulas, contexts and conclusions are unchanged.
\end{theorem}
\begin{proof}
For $\Ncal$ to $\Gcal$, induct on the native proof.
Hyp, $I_0$, $I_S$, QE and DE are common. Translate native transport by
Theorem~\ref{new:transport}, DI-case by Lemma~\ref{new:dcase}, and
QI-progress by Lemma~\ref{new:qprog}.

For $\Gcal$ to $\Ncal$, translate G by Lemma~\ref{new:gsound}.
A case-mode DI is native DI-case. For a progress-mode DI, the witness
is zero and the matrix is an equation. Native QI-progress applied to
$B(0,y)$ gives $\Q_yB(0,y)$. Under $B(0,0)$, QE gives $B(0,Sz)$;
native DI-case with witness zero then gives $\D_{x,y}B(x,y)$.
The remaining rules are identical.

All constructions retain arbitrary $\Gamma$ and use finitely many fresh
parameters. Syntax tests are decidable metatheoretic operations on
constructors; they do not invoke object-theoretic excluded middle.
\end{proof}

\begin{corollary}[Numerical substitution]\label{new:subst}
Capture-free substitution of numerical terms into derivations is an
admissible metarule of $\Gcal$.
\end{corollary}
\begin{proof}
Translate to $\Ncal$, apply its substitution theorem, and translate back.
\end{proof}

This is a derived transformation of proofs, not a claim that substituting
into one primitive G or DI node always preserves that primitive instance.

\subsection{The precise iPRA target}

Here iPRA means full first-order intuitionistic logic with a symbol for
each fixed primitive-recursive definition, its defining equations,
successor injectivity, successor/zero separation, and induction for
quantifier-free formulas with numerical parameters and retained
assumptions. The target is the \code{PRAAx} theory in the pinned
\code{PR.lean}~\cite{native}. Arbitrary-formula induction is not added.

There is no internal universal primitive-recursive evaluator.
Translations recurse externally on each fixed finite description.
An iterator update $h(u,\vec p)$ is elaborated to a function symbol
$i_h(n,a,\vec p)$ satisfying
\[
 i_h(0,a,\vec p)=a,\qquad
 i_h(Sn,a,\vec p)=h(i_h(n,a,\vec p),\vec p).
\]
This is primitive recursion whose update ignores the recursion counter.

In the reverse direction, the inherited scalar arithmetic construction
first derives predecessor, then actual-equality reflection and
quantifier-free induction, and then pairing and the compiler for every
fixed primitive recursion. For a target recursion with base $g$ and
step $h$, its state records the counter and current value:
\begin{align*}
 V(n,\vec p)&=
 I\bigl(n,\pi(0,g^*(\vec p));
 u.\pi(SL(u),h^*(L(u),H(u),\vec p))\bigr),\\
 f^*(n,\vec p)&=H(V(n,\vec p)).
\end{align*}
Here $\pi,L,H$ are the internally proved pairing and projections.
Equation induction proves $L(V(n,\vec p))=n$ and the target defining
equations. The inherited appendices reproduce the dependency chain;
pairing is not assumed to bootstrap predecessor.

\subsection{Ordinary logical connectives are definable}

The inherited formula selector $\operatorname{sel}(v,A,B)$ equals $A$
at zero and $B$ at a successor. The gate $\operatorname{gate}(v,A)$ is
trivial at zero and equivalent to $A$ at a successor. Their definitions
recurse on the native formula syntax; the appendices supply the
construction and proofs. With the parity iterator $\operatorname{par}$,
one can define
\begin{align*}
 \bot&:=0=S0,\\
 A\to B&:=\Q_y\operatorname{sel}(y,A,B),\\
 A\land B&:=\Q_y\operatorname{gate}
      (y,\operatorname{sel}(\operatorname{par}(y),A,B)),\\
 \exists x\,A(x)&:=\D_{x,y}\operatorname{gate}(y,A(x)),\\
 A\lor B&:=\exists v\,\operatorname{sel}(v,A,B),\\
 \forall x\,A(x)&:=A(0)\land\Q_y A(y).
\end{align*}
All newly bound names are fresh. For the selector and gate construction on a formula containing $\Q$,
first use its proved representation by $\D$ with an unused witness
coordinate; the appended Eq/D construction then applies. These are
definitions with derived introduction and elimination rules, not
uncounted logical axioms.

In an ordinary first-order realization one must explicitly retain the
bridge $(0=S0)\to\bot$ to genuine falsity. Positive equations alone admit
a singleton structure. The appended ordinary-bridge proof includes this
sentence and numerical coverage; neither is silently inferred from a
purely positive equational theory.

\subsection{The two translations and their composites}

Let $F$ be the native-to-iPRA translation: compile every iterator
description to its primitive-recursive function and read $\Q,\D$ by the
ordinary formulas in Section~\ref{new:rules}. Let $U$ be the
iPRA-to-native translation: compile each primitive-recursive description
to a scalar iterator term and each logical connective to its native
macro. Both domains are the whole numerical domain, defined by $x=x$,
and both equality interpretations are actual equality.

\begin{theorem}[Bi-interpretability and equiconsistency]\label{new:biinterp}
The full theory presented by $\Gcal$ is bi-interpretable and
equiconsistent with iPRA. For every native $A$ and ordinary iPRA formula
$\varphi$,
\[
 A\dashv\vdash_{\Gcal}U(F(A)),
 \qquad
 \ipra\vdash\varphi\leftrightarrow F(U(\varphi)).
\]
Both numerical composite isomorphisms have graph $x=y$. In particular,
\[
 \Gcal\vdash 0=S0
       \quad\Longleftrightarrow\quad
 \ipra\vdash\bot.
\]
\end{theorem}
\begin{proof}
The inherited native/iPRA translations preserve all sequents and have
these round trips. The formal source statements are
\code{native_iff_pra}, \code{pra_iff_native},
\code{native_roundTrip}, \code{pra_native_roundTrip}, and
\code{native_inconsistent_iff} in \code{NativeBiInterpretation.lean}.
Their arithmetic and logical constructions are reproduced in the
appendices.

For completeness, the arithmetic composite is not justified merely by
agreement on standard natural numbers. Equation induction proves
uniqueness for the two defining equations of each iterator or primitive
recursion, keeping all parameters fixed. External induction on the
finite description shows that its double translation satisfies the same
defining equations; internal uniqueness proves equality at every input.
The identity graph is therefore internally total, bijective and
operation preserving. Structural induction on formulas, using the
proved meanings of the logical macros, gives both formula composites.

Theorem~\ref{new:main} transfers every native sequent in these arguments
to $\Gcal$, without changing any term, formula or domain. Hence it
transfers both interpretation composites and the contradiction
equivalence, not only mutual interpretability.
\end{proof}

\section{Executable artifacts and size accounting}
\label{new:implementation}

\label{implementation:strict}

The delivered \code{code/global_rewrite.py} is a complete standalone
checker for the seven-rule calculus. Its input is an exact four-key JSON
object containing a nonnegative integer \code{variables}, a list
\code{context} of retained formulas, a formula \code{goal}, and a
primitive \code{proof} tree. It validates the full tagged Eq/Q/D language
and numerical variable scopes. The input contract includes open sequents
with the declared numerical parameters, as well as closed theorems.
Acceptance requires a normal exit and exactly \code{ACCEPT} followed by
one newline. The program imports only Python's standard JSON and system
modules. Neither a producer nor a theorem database is trusted.

Numerical terms are encoded as \code{['v',i]}, \code{['z]},
\code{['s',a]} and \code{['i',n,a,h]}; the final argument of the iterator
binds one variable. Formulas are \code{['=',s,t]}, \code{['q',B]} and
\code{['d',B]}. Variables use de Bruijn indices, with index zero the most
recent binder. The body of D has its case variable at index zero and
its witness variable at index one. The RW instruction contains only its
two premise proofs: \code{['rw',p,q]}. It contains no numerical or formula
context to be instantiated.

\subsection{One traversal with three precise modes}

The helper \code{walk(a,t,s,k,d)} uses the same structural recursion for
lifting, substitution and global replacement. Here $d$ records the number
of binders traversed. When a source pattern is supplied, a node matching
the source lifted through $d$ binders is replaced by the target lifted
through those binders. The traversal stops at that match. An unmatched
variable below depth $d$ remains bound; an unmatched index at least $d$
is changed by $k$. The supported calls are:

\begin{center}\small
\begin{tabular}{@{}lll@{}}
\toprule
Operation & Source pattern & Index offset\\
\midrule
Lifting by $k$ & none & $k$\\
Substitution, dropping one variable & variable zero & $-1$\\
Successor instance under a retained eigenvariable & variable zero & $0$\\
Global replacement & the proved source term & $0$\\
\bottomrule
\end{tabular}
\end{center}

The matching test precedes descent into children, which enforces the stated
top-down replacement convention. In particular, newly inserted material
is not visited again. The pattern and image are ordinary checked terms;
the internal value \code{None} denotes the absence of a pattern and is
not numerical syntax.

For DI, the matrix is first instantiated at the supplied witness. The
progress hypothesis is retained exactly when the matrix is an equation
and the witness is the literal zero term; otherwise the retained
hypothesis is its zero case. In each branch, the ambient context is lifted
when the eigenvariable is introduced. DE similarly lifts the context and
the claimed conclusion before checking its continuation. This implements
freshness through the scope discipline rather than by a convention imposed
on the certificate producer.

\subsection{Whole-program measurements}

The original X comparator, \code{reference/one_context_seven.py}, is
\emph{our matched seven-rule adaptation} of the incoming
\code{one_context/direct} checker: its two introduction constructors are
replaced by the already established witness-sensitive W introduction.
It is not represented as an unchanged published 3,181-byte file. The
callback and unified-traversal controls retain that X rule and change only
the traversal. Relative to the unified X control, the new program changes
only the equality-rule branch.

All rows charge an entire standalone module, including syntax validation,
scope checks, inference checking, JSON input and the verdict driver.
Actual docstrings are removed, and CPython 3.12.14
\code{ast.unparse} plus one final newline gives the normalized source.
Significant-token counts omit comments, encoding and whitespace markers;
AST counts include every node of \code{ast.walk}. Gzip uses level 9,
\code{mtime=0}, and zlib 1.3.2, including the gzip wrapper.

\begin{center}\small
\begin{tabular}{@{}p{0.46\linewidth}rrrr@{}}
\toprule
Complete standalone module & Bytes & Tokens & AST & Gzip-9\\
\midrule
One-context X: original traversal & 3181 & 1029 & 1089 & 976\\
One-context X: callback traversal & 3215 & 1059 & 1124 & 983\\
One-context X: unified traversal & 3206 & 1044 & 1096 & 987\\
Template-free global replacement & 3147 & 1027 & 1072 & 983\\
\bottomrule
\end{tabular}

\end{center}

The new module is 34 normalized bytes, two significant tokens and
17 AST nodes smaller than the original-traversal X control, a byte
reduction of approximately $1.07\%$. Its gzip size is seven bytes larger;
compression is reported as supplementary information. With exactly the
same unified traversal on both sides, removing the supplied context and
its checking work saves 59 bytes, 17 tokens and 24 AST nodes.

These are modest matched implementation improvements. The compact
arity-code frontier uses a different certificate language and input
contract. The table does not establish a new corpus-wide minimum, and
does not provide a lower bound on possible checkers. The mathematical
contribution is the elimination of the context parameter, together with
the large-marker completeness proof; code size is measured separately.

\subsection{Executed checks and their limits}

The command \code{python3 tests/verify_strict.py} regenerates all
certificates and the measurement table. The authoritative receipts are
\code{evidence/strict_validation.json} and
\code{evidence/strict_metrics.json}; the latter records source hashes,
raw and normalized byte counts, and exact runtime versions.

There are 38 named fixtures: 14 valid derivations and 24 expected
rejections. Four standalone programs, each run under ordinary Python and
\code{python -O}, give 304 process checks, with zero mismatches. The X
controls receive explicit X translations of the R proofs. Their retained
contexts, variable counts and claimed conclusions are identical. A
separate fixed-seed audit compares 600 liftings, 200 ordinary
substitutions and 200 substitutions without dropping the eigenvariable
against the original implementation, and checks 200 root replacements:
1,200 further comparisons, all matching.

Positive cases include large-marker congruence with accidental passive
occurrences of the source, simultaneous replacement in both equation
endpoints, nested iterator binders, a closed congruence theorem, and
equation induction that uses its hypothesis. The independent logical
naming construction exercises DI, DE and QE as well. Rejection controls
cover bound-variable capture, the wrong top-down convention, a case
hypothesis substituted for a required progress hypothesis, incorrect
computation endpoints, wrong premise kinds, malformed syntax, Boolean
and negative indices, and escaping witnesses.

During review, a preliminary callback traversal failed to preserve an
unmatched variable leaf: it attempted to descend into the variable's
integer payload. The nested-binder fixtures detected the defect. That
prototype and its preliminary size are not the delivered result. The
unified traversal handles every variable explicitly, and the final
receipts refer to the corrected source hashes.

The implementation audit also checks the inference invariants directly:
each checked premise has a well-scoped formula; lifting preserves the
bound/free distinction; substitution removes exactly its designated
variable when requested; and replacement inserts a term only after
lifting it through the encountered binders. These are written source
arguments supported by executable cases. The finite tests are not a
proof-assistant verification of the Python implementation or of
bi-interpretability. As for the baseline, the universal proof-system
statement abstracts from finite process memory and recursion limits;
a resource failure is never acceptance.

\subsection{A compact forward supplement}
\label{supp:compact-forward}

The directory \code{compact/} supplies a second implementation, using the
inherited raw arity representation and a synthesized conclusion. Each
non-hypothesis proof is a pair consisting of a major proof and rule data.
When the major is $s=t$, its data are only a minor proof $u=v$; the program
returns $v=\Rep{s}{t}{u}$. No equality template is present in that packet.
This supplement is distinct from the strictly parsed tagged interface
measured above: administrative aliases retain their inherited meanings,
and a strict public-goal grammar validator is not included.

The compact arithmetic packets omit literal-zero initializers. The new
congruence compiler therefore uses the uniform aliases
\[
 C_s=I(2,1;u.s),\qquad K=I(0,C_s;u.N),
\]
where the closed numeral $N$ is larger in syntax-tree size than the desired
context. An allowed successor computation gives $C_s=s$, and an allowed
base computation gives $K=C_s$. Root transitivity gives $K=s$, including
for the literal source $s=0$. The oversized-alias proof then applies
without presupposing missing computation instances or existential naming.

\begin{center}\small
\begin{tabular}{@{}lrrrr@{}}
\toprule
Complete forward module & Bytes & Tokens & AST & Gzip-9\\
\midrule
Global replacement &865&397&464&437\\
Crossed-rule control &830&384&448&414\\
\bottomrule
\end{tabular}
\end{center}
These are whole-module normalized measurements under CPython 3.12.14 and
zlib 1.3.2. Exact delivered sources contain 862 and 827 bytes and compress
to 435 and 411 bytes respectively. The gzip files preserve those exact
sources; \code{compact/metrics.json} records both raw and normalized
identities. The two forward programs use different packet layouts, so
this pair is not a matched ablation of only the equality rule. Neither
program improves the 520-byte compact record.

\code{compact/verify.py} records 25 global-replacement cases in normal
and optimized Python, giving 50 process expectations, plus 17 crossed-rule
control cases in both modes, giving 34 more. Two deliberate-mutant runs
confirm the supplied closed-false-proof attack on the shared-sentinel
defect; the corrected program rejects it. All expectations pass. Positive
tests include accidental passive occurrences, both dependent equation
endpoints, literal-zero sources, two nested iterator binders, and a DE
continuation that actually uses its fresh witness before discharging it.
Negative tests cover witness escape, future names, and a non-unary DI goal.

The implementation keeps the replacement-source sentinel distinct from
the missing-image sentinel, makes the new source argument keyword-only,
and restores the old scope floor before checking a DE result. The README
documents these guards and the exact input contract. The scripts run from
any working directory with the Python standard library alone. This is
finite-tested supplementary code with a written mathematical explanation;
no proof-assistant verification of its implementation is claimed.

\subsection{Current comparison programs and their input contracts}

The final comparison boundary is commit
\code{2a724ea8ec6665f91aa74c892414584b36ca698d}, whose recorded commit time is
2026-11-12 01:58:39 UTC. After the earlier filed snapshot, the last check
adds the \code{checked_refinement}, \code{validated_core} and
\code{one_sided_transport} archives. The first two explicitly refine
implementations of existing mathematical theories. The last retains a
supplied numerical template in its crossed equality rule. None of these
three reports supplies the present global replacement primitive and
oversized-marker proof. This bounds the comparison to inspected material.

The current compact coordinate envelope is
\[
            256\text{ tokens}\;/\;278\text{ AST nodes}\;/\;520\text{ bytes}.
\]
Three different programs attain these coordinates. In particular, the
520-byte program has 263 tokens and 303 AST nodes; it replaces the earlier
521-byte profile. Its one-byte improvement increases that profile's AST
count from 300 to 303. The complete sources and original package claims
are copied unchanged under \code{reference/current_compact/}.

\begin{center}
\small
\begin{tabular}{lrrrr}
\toprule
Comparison program & Tokens & AST nodes & Bytes & Gzip-9\\
\midrule
Compact byte profile & 263 & 303 & 520 & 329\\
Compact token profile & 256 & 283 & 594 & 346\\
Compact AST profile & 260 & 278 & 602 & 344\\
Native Eq/D, goal checked & 358 & 413 & 788 & 460\\
Guarded Eq/D, goal checked & 407 & 459 & 921 & 507\\
Native Eq/D, fully typed & 390 & 454 & 878 & 498\\
Native Eq/Q/D, fully typed & 401 & 468 & 897 & 505\\
\bottomrule
\end{tabular}
\end{center}

The compact profiles assume a canonical guarded public goal and omit an
included public-goal validator. The 788-byte native profile validates
the full Eq/D public grammar, with Q available internally, but does not
require exact integer/list typing of every certificate field. The
fully typed profiles impose that stronger contract; the 897-byte profile
also admits every Eq/Q/D public formula. These distinctions prevent a
raw byte comparison from being mistaken for a comparison of identical
interfaces. The present tagged and compact implementations retain the
contracts stated in their own preceding subsections.

Every displayed source was remeasured locally with CPython 3.12.14 using
the whole module after \code{ast.unparse(ast.parse(source))} plus one LF.
The token count excludes comments, layout, encoding and end markers;
all AST nodes are counted. Gzip uses level 9, timestamp zero and the
complete wrapper, and is informational. These measurements reproduce
the imported packages' recorded values.

The three new ZIPs were authenticated against the Git blob identities
returned by the repository API. Each copied member was then checked
against its ZIP bytes. Archive identities, member paths, SHA-256 hashes,
individual input contracts and reproduced metrics are recorded in
\code{evidence/comparison_sources.json}; original claims and metric files
are retained beside the sources. The imported packages' entire process
test suites were not rerun for this size comparison.

The one-template seven-rule matched control used earlier is a derived
control, not an unchanged upstream file. Its exact ancestor is the
archived \code{direct/one_context_checker.py}. Replacing its separate
case/progress introduction branches by the existing witness-sensitive
introduction produces \code{reference/one_context_seven.py}, with no
other change. The exact difference is included as
\code{evidence/one_context_seven_origin.diff}. Further shared-walker
variants are identified separately in the matched implementation audit.


\section{Verification status, limits and reproduction}
\label{new:verification}

The new results are universal written proofs, independently reviewed
within this investigation and accompanied by executable proof
certificates. They have not been checked in Lean or Rocq. The inherited
Lean sources were inspected, but the complete upstream project was not
rebuilt here. Existing Lean theorems belong to the inherited bridge;
they are not presented as a formalization of G. The copied Lean files
are reference excerpts, not a standalone Lake project.

The checkers are standalone and do not import the proof producers or a
theorem database. Positive tests exercise genuine primitive proof
packets; negative controls exercise binder capture, wrong occurrence
semantics, malformed syntax and rule-scope failures. Finite tests do not
replace the universal mathematical proof.

Implementation traps were found and fixed during checking.
A generic tree walker must preserve an unmatched variable leaf instead
of recursing into its integer payload. In the compact shared opener,
the missing-image sentinel must differ from the replacement-source
sentinel, and the additional source argument must be keyword-only;
otherwise existing arity guards can be bypassed. Its DE test must
actually use the opened witness to exercise allocation order. The
archived implementations and measurements are the corrected versions.

The main result is therefore a concrete seven-rule alternative with
a strictly smaller equality-certificate interface, complete derivability
equivalence, explicit interpretation composites, and reproducible
implementations. The outstanding optimization problem is whether this
template elimination can also improve the smallest complete checker
under a fixed compact input contract. No absolute optimum or new
compact record is claimed.

Run \code{python3 reproduce.py} from the extracted directory to regenerate
measurements and process receipts. Run \code{python3 build.py} to flatten
the attributed appendices and compile the PDF using TeX Live.
The archive includes the editable modular source, a standalone
\code{ipra_template_free.tex}, its PDF, code, proof producers, fixtures,
comparison sources, source hashes and the verification record.

\clearpage
\appendix
\section*{Inherited foundations: attribution and dependency boundary}
\addcontentsline{toc}{section}{Inherited foundations}

The following five sections are reproduced from the pinned
\code{Research/certified_zero/reference/} sources under ProveIt's MIT No
Attribution license. Their original wording and historical attributions
are retained. References to a new result in these reproduced sections
refer to their older packages, not to novelty in the present report.

Some appended proofs use reusable clauses and a native Eq/D presentation.
A closure clause $\langle B(0);y.B(y)\rangle$ has the meaning of
$\Q_yB(y)$; the ordinary-bridge appendix explicitly proves the
correspondence, including the unused-witness representation of $\Q$.
These alternate presentations are inherited foundations. The exact
primitive language of this proposal is the Eq/Q/D language of
Section~\ref{new:rules}.

The new mathematical construction and its proof are
Sections~\ref{new:rules}--\ref{new:equivalence}.

% INHERITED_APPENDICES

\clearpage
\begin{thebibliography}{9}
\bibitem{snapshot}
Vladimir Reshetnikov et al., \emph{Concise presentations of intuitionistic
primitive recursive arithmetic}, ProveIt repository, source snapshot
\code{a4e5e96c456303cf231eab2aeeb3291d8889833f}.
\url{https://github.com/VladimirReshetnikov/ProveIt/tree/a4e5e96c456303cf231eab2aeeb3291d8889833f/Logic/PrimitiveRecursiveArithmetic/ConciseIPRA}.

\bibitem{native}
ProveIt, \emph{Native.lean, NativeOneIntro.lean, NativeBiInterpretation.lean,
BiInterpretation.lean, PR.lean}, same snapshot. Exact copied-file
identities and URLs appear in \code{evidence/sources.json}.

\bibitem{comparison}
ProveIt, \emph{one\_context, kernel\_refinement, compact\_update,
uniform\_progress}, filed snapshot
\code{521884df22e8f2dac0b7066095a42ae975ea4166};
additional incoming archives checked through
\code{2a724ea8ec6665f91aa74c892414584b36ca698d}.
\url{https://github.com/VladimirReshetnikov/ProveIt/tree/2a724ea8ec6665f91aa74c892414584b36ca698d/Logic/PrimitiveRecursiveArithmetic/ConciseIPRA}.
See the bounded source audit in the archive.

\bibitem{forward}
ProveIt, inherited \emph{Equality reflection and quantifier-free induction},
\path{Research/certified_zero/reference/equality_reflection.tex}, source
snapshot. Historical forward-left attribution retained in the appendix.

\bibitem{erasure}
ProveIt, inherited \emph{Scalar predecessor bootstrap},
\path{Research/certified_zero/reference/scalar_predecessor.tex}, source snapshot.

\bibitem{progress}
ProveIt, inherited \emph{Primitive recursion from scalar iteration},
\path{Research/certified_zero/reference/primitive_recursion.tex}, source snapshot.
\end{thebibliography}
\end{document}
